\section{The quotient-frame orbit}

Section~4 constructed the natural quotient frame

\[
\mathcal P_{0}
\]

and its deck involution

\[
\tau_{0}.
\]

Section~6 established the direct-product decomposition

\[
A=\Aut(X)=K\times C,
\]

where

\[
K\cong S_{5}
\]

and

\[
C\cong S_{3}.
\]

We now interpret the three involutions of the complementary factor \(C\)
as three conjugate regular quotient frames.

\subsection{The natural deck involution inside \(C\)}

% STATUS: Certified computational alignment
% SOURCE: Project 42 explicit quotient-frame certificate 004
% COMMIT: 1604b16

The covering-space and homogeneous-space descriptions of \(X\) initially
produce their permutation groups in different coordinate systems. The
explicit labeled identification used in Section~5 aligns these two
descriptions.

\begin{proposition}
\label{prop:natural-deck-in-centralizer}

Under the certified labeled identification of the two graph models, the
natural deck involution satisfies

\[
\tau_{0}\in C.
\]

In particular, \(\tau_{0}\) is one of the three involutions of the
complementary group

\[
C\cong S_{3}.
\]

\end{proposition}

\begin{proof}

The explicit permutation certificate transports the natural deck
involution from the canonical-cover model to the homogeneous-space model
and verifies that it centralizes the transitive subgroup \(K\).

Thus

\[
\tau_{0}\in C_{\Aut(X)}(K)=C.
\]

It is nonidentity and has order two. Hence, under the identification
\(C\cong S_{3}\), it is a transposition.
\end{proof}

This is the only coordinate-alignment input required in this section.

\subsection{The three involutions}

Let

\[
\tau_{0},
\qquad
\tau_{1},
\qquad
\tau_{2}
\]

denote the three involutions of \(C\).

All three are conjugate within \(C\). Since \(\tau_{0}\) is
fixed-point-free on \(V(X)\), each \(\tau_i\) is fixed-point-free.

For each

\[
i\in\{0,1,2\},
\]

let

\[
\mathcal P_i
\]

be the orbit partition of \(\langle\tau_i\rangle\) on \(V(X)\).

Each partition therefore consists of fifteen unordered pairs.

The conjugation action of \(C\) on its three involutions induces the
following action on the three partitions:

\[
\rho(\tau_{0})=(1\ 2),
\]

\[
\rho(\tau_{1})=(0\ 2),
\]

and

\[
\rho(\tau_{2})=(0\ 1).
\]

Equivalently, in image-list notation,

\[
\rho(\tau_{0})=[0,2,1],
\qquad
\rho(\tau_{1})=[2,1,0],
\qquad
\rho(\tau_{2})=[1,0,2].
\]

\subsection{The three pair partitions}

The natural frame \(\mathcal P_0\) is the fiber partition constructed in
Section~4. The other two frames are obtained by conjugating
\(\mathcal P_0\) within \(C\).

More explicitly, if

\[
c\tau_{0}c^{-1}=\tau_i,
\]

then

\[
c(\mathcal P_0)=\mathcal P_i.
\]

Thus the three pair partitions are intrinsic to the action of the
complementary \(S_3\), although the label \(0\) records the frame selected
by the original covering construction.

The complete labeled pair lists and the associated involutions are
recorded in Appendix~B.

\subsection{The quotient graphs}

% STATUS: Conceptual deduction from Proposition 7.1

\begin{proposition}
\label{prop:three-frame-quotients}

For each

\[
i\in\{0,1,2\},
\]

the partition \(\mathcal P_i\) is a regular two-fold quotient frame and

\[
X/\langle\tau_i\rangle
\cong
L(P).
\]

\end{proposition}

\begin{proof}

The assertion is established for \(i=0\) in Section~4.

For arbitrary \(i\), choose \(c\in C\) such that

\[
c\tau_{0}c^{-1}=\tau_i.
\]

Conjugation by the graph automorphism \(c\) carries the orbit partition of
\(\tau_0\) to the orbit partition of \(\tau_i\). It also transports the
locally bijective quotient projection for \(\mathcal P_0\) to a locally
bijective quotient projection for \(\mathcal P_i\).

Therefore the quotient graphs are isomorphic:

\[
X/\langle\tau_i\rangle
\cong
X/\langle\tau_0\rangle
\cong
L(P).
\]
\end{proof}

\subsection{Frame stabilizers}

% STATUS: Conceptual deduction from the direct-product theorem

\begin{proposition}
\label{prop:three-frame-stabilizers}

For each

\[
i\in\{0,1,2\},
\]

the stabilizer of \(\mathcal P_i\) in \(A\) is

\[
A_{\mathcal P_i}
=
K\times\langle\tau_i\rangle
\cong
S_{5}\times C_{2}.
\]

\end{proposition}

\begin{proof}

The partition \(\mathcal P_i\) determines its block-swap involution
\(\tau_i\) uniquely. Hence an automorphism stabilizes \(\mathcal P_i\) if
and only if it centralizes \(\tau_i\).

Since

\[
A=K\times C
\]

and \(K\) centralizes all of \(C\),

\[
C_A(\tau_i)
=
K\times C_C(\tau_i).
\]

The centralizer of a transposition in \(S_3\) is the subgroup generated by
that transposition. Therefore

\[
C_C(\tau_i)
=
\langle\tau_i\rangle,
\]

and the result follows.
\end{proof}

In particular,

\[
|A_{\mathcal P_i}|=240
\]

for every \(i\).

\subsection{The forty-five-pair census}

% STATUS: Conceptual deduction

The three partitions do not share blocks.

\begin{proposition}
\label{prop:forty-five-pairs}

The partitions

\[
\mathcal P_0,
\qquad
\mathcal P_1,
\qquad
\mathcal P_2
\]

are pairwise block-disjoint. Consequently,

\[
\left|
\mathcal P_0
\cup
\mathcal P_1
\cup
\mathcal P_2
\right|
=
45.
\]

\end{proposition}

\begin{proof}

First observe that the centralizer \(C\) acts semiregularly on \(V(X)\).
Indeed, suppose \(c\in C\) fixes a vertex \(v\). Since \(K\) is transitive,
every vertex has the form \(v^k\) for some \(k\in K\). Commutation gives

\[
(v^k)^c
=
(v^c)^k
=
v^k.
\]

Thus \(c\) fixes every vertex, and therefore \(c=1\).

Now suppose that distinct partitions \(\mathcal P_i\) and
\(\mathcal P_j\) shared a block

\[
\{v,w\}.
\]

Then

\[
\tau_i(v)=w=\tau_j(v),
\]

and hence the nonidentity element

\[
\tau_j\tau_i
\]

would fix \(v\). This contradicts semiregularity of \(C\).

Therefore the three partitions are pairwise block-disjoint. Each contains
fifteen blocks, so their union contains

\[
3\cdot15=45
\]

distinct unordered pairs.
\end{proof}

\subsection{The quotient-frame orbit}

% STATUS: Conceptual theorem after one certified alignment input

\begin{theorem}
\label{thm:quotient-frame-orbit}

The natural quotient frame has a three-member automorphism orbit:

\[
\frorbit
=
A\cdot\mathcal P_0
=
\{
\mathcal P_0,
\mathcal P_1,
\mathcal P_2
\}.
\]

Every member of this orbit is a regular two-fold quotient frame with
quotient isomorphic to \(L(P)\). The induced action of \(A\) on
\(\frorbit\) has image \(S_3\), and

\[
A_{\mathcal P_i}
=
K\times\langle\tau_i\rangle
\cong
S_5\times C_2
\]

for each \(i\).

Moreover, the three frames are pairwise block-disjoint and together
contain forty-five distinct unordered vertex pairs.
\end{theorem}

\begin{proof}

The group \(K\) centralizes \(C\), so it fixes each of the three
involutions \(\tau_i\) under conjugation and therefore fixes each
partition \(\mathcal P_i\).

The group \(C\cong S_3\) acts transitively by conjugation on its three
transpositions. Hence it acts transitively on

\[
\{
\mathcal P_0,
\mathcal P_1,
\mathcal P_2
\}.
\]

The conjugation action of \(S_3\) on its transpositions is faithful, so the
induced action on the three frames has image \(S_3\).

The quotient, stabilizer, and block-disjointness statements follow from
Propositions~\ref{prop:three-frame-quotients},
\ref{prop:three-frame-stabilizers}, and
\ref{prop:forty-five-pairs}.
\end{proof}

No assertion is made that every admissible regular \(L(P)\)-quotient frame
belongs to this orbit.
